\begin{table}[H]
\centering
\caption{NFR interpolation versus exact assignment at matched constraints}
\label{tab:stable-update}
\small
\begin{tabular}{rrrrrrrr}
\toprule
 & Queue & \multicolumn{3}{c}{NFR (\%)} & \multicolumn{3}{c}{Accuracy (\%)} \\
\cmidrule(lr){3-5}\cmidrule(lr){6-8}
$\alpha$ & shift (\%) & Interp. & QS & QS+flips & Interp. & QS & QS+flips \\
\midrule
0.0 & 0.00 & 0.00 & 12.28 & 0.00 & 52.92 & 60.56 & 52.92 \\
0.2 & 1.53 & 2.30 & 12.38 & 2.29 & 56.08 & 60.80 & 57.70 \\
0.5 & 3.28 & 6.49 & 12.56 & 6.47 & 59.65 & 61.02 & 60.21 \\
0.8 & 5.67 & 10.79 & 12.98 & 10.76 & 60.99 & 61.14 & 61.04 \\
1.0 & 7.07 & 13.28 & 13.28 & 13.28 & 61.16 & 61.16 & 61.16 \\
\bottomrule
\end{tabular}
\begin{minipage}{0.94\linewidth}
\footnotesize\emph{Note:} The data-generating process is fixed across release and planning samples. The candidate observes an additional independent signal and is deployed only after beating the incumbent on an independent out-of-sample release set. Each row compares probability interpolation (Interp.) with exact assignment under the candidate posterior at the interpolation point's queue shift (QS), and at both its queue shift and its expected negative flips (QS+flips). Results average 1,000 planning batches from 200 accepted updates. Accuracy and NFR use labels revealed only after every assignment is fixed.
\end{minipage}
\end{table}
